\chapter{Circuits of Kraus families on a geometry}
\label{ch:circuits}

This chapter uses the notation of Chapter~\ref{ch:sites}: $V$ is a finite set of
sites, a system is a family $X=(X_v)_{v\in V}$ of finite basis sets, and
$\mathcal H_X=\bigotimes_v\C^{X_v}$. A circuit acts layer by layer by finite Kraus
families on disjoint cells of a geometry, and its basis sets may change from layer
to layer, so that ancillas can be prepared and discarded and local dimensions can
change. Later layers may depend on the outcomes of earlier ones. Unitary circuits,
channels, measurement with classical feedback, block encodings and circuits with
oracle queries are circuits of this one kind, distinguished by properties of their
Kraus families \cite{NielsenChuang2010}. The chapter proves, once for all of them,
the composition of circuits, the subadditivity of depth and gate count, trace
preservation, and the backward light cone.

\section{Geometries, gates and layers}
\label{sec:circuits_layers}

\begin{definition}[Geometry]
    \label{def:circuits_geometry}
    \lean{QuantumCircuit.Geometry, QuantumCircuit.Geometry.neighbourhood,
        QuantumCircuit.Geometry.subset_neighbourhood,
        QuantumCircuit.Geometry.restrict}
    \leanok
    A geometry on $V$ is a set $\mathcal C$ of cells, each with a region
    $C_b\subseteq V$ and a cost $c_b\in\N$. The bonds of a ring, the windows of a
    chain, the edges of a graph, the wires of one site at cost zero and arbitrary
    sets of parties are geometries. When there are finitely many cells, the
    neighbourhood $\mathcal N(R)$ of a region $R$ is $R$ together with every cell
    meeting $R$, so $R\subseteq\mathcal N(R)$.
    The cells contained in a region $D$ form a geometry on $D$.
\end{definition}

\begin{lemma}[Monotonicity of neighbourhoods]
    \label{lem:circuits_neighbourhood_mono}
    \lean{QuantumCircuit.Geometry.neighbourhood_mono,
        QuantumCircuit.Geometry.iterate_neighbourhood_mono}
    \leanok
    \uses{def:circuits_geometry}
    For a geometry with finitely many cells, if $R\subseteq R'$ then
    $\mathcal N(R)\subseteq\mathcal N(R')$, and if moreover
    $m\leq n$ then $\mathcal N^m(R)\subseteq\mathcal N^n(R')$.
\end{lemma}
\begin{proof}\leanok
    A site of $\mathcal N(R)$ lies in $R\subseteq R'$ or in a cell meeting $R$, which
    then meets $R'$. The second statement follows by induction on $n$: if
    $m\leq n$ then
    $\mathcal N^m(R)\subseteq\mathcal N^n(R')\subseteq\mathcal N^{n+1}(R')$, and
    if $m=n+1$ then $\mathcal N^{n+1}(R)=\mathcal N(\mathcal N^n(R))\subseteq
    \mathcal N(\mathcal N^n(R'))$ by the first statement.
\end{proof}

\begin{definition}[Gate families and interpretations]
    \label{def:circuits_gate_family}
    \lean{QuantumCircuit.GateFamily, QuantumCircuit.Interpretation,
        QuantumCircuit.KrausGate, QuantumCircuit.KrausGate.ev}
    \leanok
    \uses{def:kraus_family, def:sites_config}
    A gate family assigns to every region $D$ and systems $X,Y$ a set of gates on
    $D$ from $X$ to $Y$, each with a finite set of outcomes. An interpretation
    assigns to every outcome of every gate a matrix in
    $\C^{Y_D\times X_D}$. The plain gates on $D$ are the finite Kraus families
    from $\C^{X_D}$ to $\C^{Y_D}$, each interpreted as itself. Keeping the
    matrices out of the gates allows one circuit to be evaluated against
    different oracles, and allows gates that are themselves formal expressions,
    such as weighted sums of circuits.
\end{definition}

\begin{definition}[Layer]
    \label{def:circuits_layer}
    \lean{QuantumCircuit.Layer, QuantumCircuit.Layer.Free,
        QuantumCircuit.Layer.Index, QuantumCircuit.Layer.kraus,
        QuantumCircuit.Layer.krausFamily, QuantumCircuit.Layer.Forall,
        QuantumCircuit.Layer.lightCone}
    \leanok
    \uses{def:circuits_geometry, def:circuits_gate_family}
    A layer from $X$ to $Y$ consists of a set of pairwise disjoint cells $b$, each
    carrying a gate $g_b$ on $C_b$, and of a finite Kraus family
    $(k_{s,\omega})_\omega$ from $\C^{X_s}$ to $\C^{Y_s}$ at every site $s$ met by
    no cell. Its outcomes are the choices $i$ of an outcome for every cell and
    every remaining site, and its Kraus operators are
    \begin{align}
        (K_i)_{yx}&=\prod_b g_b(i_b)_{y|_{C_b},x|_{C_b}}\prod_sk_{s,i_s}(y_s,x_s).
        \notag
    \end{align}
    A property of Kraus families holds locally on the layer when it holds for
    every gate, as interpreted, and every onsite family. The backward light cone
    of a region $R$ through the layer is $R$ together with every cell of the layer
    that meets $R$.
\end{definition}

\begin{lemma}[A single gate is its placement]
    \label{lem:circuits_layer_single}
    \lean{QuantumCircuit.Layer.single, QuantumCircuit.Layer.free_single_iff,
        QuantumCircuit.Layer.kraus_single}
    \leanok
    \uses{def:circuits_layer, def:sites_place}
    The layer with one plain gate $V$ on a cell $b$ and the identity at every other
    site has the single Kraus operator $V\otimes1$, the placement of $V$ on $C_b$.
\end{lemma}
\begin{proof}\leanok
    At a site outside $C_b$ the entry of the identity is $1$ when the two
    configurations agree there and $0$ otherwise; the product of these entries is
    the condition in the formula for the entries of a placed matrix.
\end{proof}

\section{The Heisenberg picture of a layer}
\label{sec:circuits_layer_dual}

The cells of a layer and its remaining sites are its blocks; they partition $V$,
and each Kraus operator is a product of local entries over the blocks.

\begin{lemma}[Factorization along a union of blocks]
    \label{lem:circuits_layer_factorization}
    \lean{QuantumCircuit.Layer.Block, QuantumCircuit.Layer.region,
        QuantumCircuit.Layer.BlockIndex, QuantumCircuit.Layer.indexEquiv,
        QuantumCircuit.Layer.region_disjoint, QuantumCircuit.Layer.blockOf,
        QuantumCircuit.Layer.mem_region_blockOf, QuantumCircuit.Layer.blockOf_eq_of_mem,
        QuantumCircuit.Layer.blockKraus, QuantumCircuit.Layer.kraus_apply_eq_prod,
        QuantumCircuit.Layer.blockEffect, QuantumCircuit.Layer.blockEffect_eq_one,
        QuantumCircuit.Layer.unionRegion, QuantumCircuit.Layer.region_subset_unionRegion,
        QuantumCircuit.Layer.disjoint_unionRegion, QuantumCircuit.Layer.isComplPartition,
        QuantumCircuit.Layer.nearKraus, QuantumCircuit.Layer.farKraus,
        QuantumCircuit.Layer.farKraus_apply, QuantumCircuit.Layer.nearFarEquiv,
        QuantumCircuit.Layer.kraus_eq_submatrix, QuantumCircuit.Layer.sum_farKraus,
        QuantumCircuit.Layer.dual, QuantumCircuit.Layer.dual_apply',
        QuantumCircuit.Layer.dual_place}
    \leanok
    \uses{def:circuits_layer, lem:sites_partition, def:sites_place}
    Let $D$ be the union of a set $N$ of blocks of a layer. In the cut along $D$,
    every Kraus operator is $K_i=K^{\mathrm{near}}_{i_N}\otimes K^{\mathrm{far}}_{i_F}$,
    the products of the local entries over the blocks in $N$ and over the
    remaining blocks $F$. If the families on the blocks in $F$ are
    trace-preserving, then for every $A_1$ on $\C^{Y_D}$
    \begin{align}
        \sum_iK_i^\dagger(A_1\otimes1)K_i
        &=\Bigl(\sum_{i_N}(K^{\mathrm{near}}_{i_N})^\dagger A_1
        K^{\mathrm{near}}_{i_N}\Bigr)\otimes1.
        \notag
    \end{align}
\end{lemma}
\begin{proof}\leanok
    Splitting the product over the blocks into the blocks in $N$ and the others
    gives the Kronecker factorization. Then
    $K_i^\dagger(A_1\otimes1)K_i=((K^{\mathrm{near}})^\dagger A_1K^{\mathrm{near}})
    \otimes((K^{\mathrm{far}})^\dagger K^{\mathrm{far}})$, and the sum over the
    outcomes splits into a sum over $i_N$ and a sum over $i_F$. By
    Lemma~\ref{lem:sites_partition} the entry of
    $\sum_{i_F}(K^{\mathrm{far}}_{i_F})^\dagger K^{\mathrm{far}}_{i_F}$ at $(w',w)$
    is the product over the blocks in $F$ of the entries of their effects at the
    restrictions of $w'$ and $w$. Each effect is the identity, so the product is
    $1$ when $w'$ and $w$ agree on every block and $0$ otherwise; since the blocks
    cover $V\setminus D$, this is the identity.
\end{proof}

\begin{theorem}[Trace preservation of a layer]
    \label{thm:circuits_layer_tp}
    \lean{QuantumCircuit.Layer.dual_one, QuantumCircuit.Layer.isTP_krausFamily}
    \leanok
    \uses{def:circuits_layer, def:kraus_family}
    A layer whose gates and onsite families are trace-preserving is
    trace-preserving.
\end{theorem}
\begin{proof}\leanok
    \uses{lem:circuits_layer_factorization}
    Apply Lemma~\ref{lem:circuits_layer_factorization} with $N$ empty and
    $A_1=1$: the union $D$ is empty, so the near factor is the identity on a
    one-point space and $\sum_iK_i^\dagger K_i=1$.
\end{proof}

\begin{theorem}[Backward light cone of a layer]
    \label{thm:circuits_layer_light_cone}
    \lean{QuantumCircuit.Layer.dual_mem_supportedOperators_lightCone}
    \leanok
    \uses{def:circuits_layer, def:kraus_family, def:sites_supported}
    If the gates and onsite families of a layer are trace-preserving and $A$ is
    supported on $R$, then $\sum_iK_i^\dagger AK_i$ is supported on $R$ together
    with the cells of the layer that meet $R$.
\end{theorem}
\begin{proof}\leanok
    \uses{lem:circuits_layer_factorization, thm:sites_supported_subalgebra}
    Let $N$ be the blocks meeting $R$ and $D$ their union. Every site of $R$ lies in
    a block meeting $R$, so $R\subseteq D$ and $A=A_1\otimes1$ for an operator
    $A_1$ on $D$. By Lemma~\ref{lem:circuits_layer_factorization} the Heisenberg
    image is the placement of an operator on $D$, hence supported on $D$. A block
    of $N$ is either a cell meeting $R$ or a site of $R$.
\end{proof}

\section{Circuits}
\label{sec:circuits_circuits}

\begin{definition}[Circuit]
    \label{def:circuits_circuit}
    \lean{QuantumCircuit.Circuit, QuantumCircuit.Circuit.Index,
        QuantumCircuit.Circuit.kraus, QuantumCircuit.Circuit.kraus_done,
        QuantumCircuit.Circuit.kraus_step, QuantumCircuit.Circuit.krausFamily,
        QuantumCircuit.Circuit.map, QuantumCircuit.Circuit.dual,
        QuantumCircuit.Circuit.dual_step, QuantumCircuit.Circuit.seq}
    \leanok
    \uses{def:circuits_layer}
    A circuit from $X$ to $Z$ is either empty, when $X=Z$, or a layer $L$ from $X$
    to some system $Y$ followed, for every outcome $i$ of $L$, by a circuit from $Y$
    to $Z$. The dependence on $i$ is classical feedback; a circuit in which no
    continuation depends on the outcome is unbranched. A branch is a choice of an
    outcome of every layer along a path, and its Kraus operator is the product of
    the Kraus operators of its layers, the first layer on the right. The map of the
    circuit is $\rho\mapsto\sum_pK_p\rho K_p^\dagger$ over all branches $p$, and its
    Heisenberg picture is $A\mapsto\sum_pK_p^\dagger AK_p$, which is the Heisenberg
    picture of the first layer applied to those of the continuations. A finite
    sequence of layers on one system, applied from the first, is an unbranched
    circuit.
\end{definition}

\begin{theorem}[Composition]
    \label{thm:circuits_composition}
    \lean{QuantumCircuit.Circuit.append, QuantumCircuit.Circuit.comp,
        QuantumCircuit.Circuit.appendIndexEquiv, QuantumCircuit.Circuit.kraus_append,
        QuantumCircuit.Circuit.compIndexEquiv, QuantumCircuit.Circuit.kraus_comp,
        QuantumCircuit.Circuit.map_comp, QuantumCircuit.Circuit.map_done,
        QuantumCircuit.Circuit.done_comp, QuantumCircuit.Circuit.comp_done}
    \leanok
    \uses{def:circuits_circuit}
    Attaching a circuit $k_p$ to every branch $p$ of a circuit $c$ gives a circuit
    whose branches are the pairs $(p,q)$ of a branch of $c$ and a branch of $k_p$,
    with Kraus operators $K^{k_p}_qK^c_p$. For a continuation independent of $p$ the
    map of the composite is the composite of the maps, and the empty circuit is the
    identity map. The empty circuit is a two-sided identity of composition.
\end{theorem}
\begin{proof}\leanok
    Induction on $c$; the branches of the composite are rearranged by the
    associativity of dependent pairs, which gives $K_{(p,q)}=K^{k_p}_qK^c_p$. For a
    continuation $c'$ independent of $p$,
    \begin{align}
        \sum_{p,q}K'_qK_p\,\rho\,(K'_qK_p)^\dagger
        &=\sum_qK'_q\Bigl(\sum_pK_p\rho K_p^\dagger\Bigr)K'^\dagger_q,
        \notag
    \end{align}
    the map of $c'$ applied to the map of $c$. The empty circuit has the single
    Kraus operator $1$; composing with it on the right changes no layer, by
    induction on $c$.
\end{proof}

\section{Costs}
\label{sec:circuits_costs}

\begin{definition}[Depth, gate count and participation]
    \label{def:circuits_costs}
    \lean{QuantumCircuit.Circuit.total, QuantumCircuit.Circuit.total_done,
        QuantumCircuit.Layer.depth, QuantumCircuit.Layer.gateCount,
        QuantumCircuit.Layer.meets, QuantumCircuit.Circuit.depth,
        QuantumCircuit.Circuit.gateCount, QuantumCircuit.Circuit.participation,
        QuantumCircuit.Layer.cellCost, QuantumCircuit.Layer.gateCount_eq_sum,
        QuantumCircuit.Layer.meets_eq_card, QuantumCircuit.Layer.depth_eq_min_one}
    \leanok
    \uses{def:circuits_circuit}
    A cost of layers determines a cost of circuits by accumulation along the most
    expensive branch. A weight of gates assigns to every gate on a cell a natural
    number, and the cell cost of a layer is the sum of the weights of its gates. The
    gate count of a layer is the cell cost for the weight $c_b$, the sum of the
    costs of its cells; the number of gates of the layer meeting a site $s$ is the
    cell cost for the weight that is $1$ on a cell of positive cost containing $s$
    and $0$ otherwise. The depth of a layer is $\min(1,m)$, where $m$ is the number
    of its cells of positive cost. Onsite families cost nothing. A circuit is of
    depth $T$ when its depth is at most $T$.
\end{definition}

\begin{theorem}[Subadditivity of costs]
    \label{thm:circuits_cost_subadditive}
    \lean{QuantumCircuit.Circuit.total_append, QuantumCircuit.Circuit.total_comp,
        QuantumCircuit.Circuit.depth_comp, QuantumCircuit.Circuit.gateCount_comp,
        QuantumCircuit.Circuit.participation_comp, QuantumCircuit.Circuit.total_mono,
        QuantumCircuit.Layer.depth_le_gateCount,
        QuantumCircuit.Circuit.depth_le_gateCount,
        QuantumCircuit.Circuit.total_seq_le, QuantumCircuit.Circuit.depth_seq_le,
        QuantumCircuit.Circuit.gateCount_seq_le}
    \leanok
    \uses{def:circuits_costs, thm:circuits_composition}
    The cost of a circuit with continuations $k_p$ attached is at most the cost of
    the circuit plus the largest cost of a continuation. In particular depth, gate
    count and the number of gates meeting a site are subadditive under
    composition. The depth of a circuit is at most its gate count. The cost of a
    sequence of layers is at most the sum of the costs of its layers.
\end{theorem}
\begin{proof}\leanok
    Induction on the circuit. The depth of a layer is at most its gate count,
    because a cell of positive cost contributes at least $1$ to the sum of costs,
    and accumulation along branches is monotone.
\end{proof}

\section{Trace preservation and the backward light cone}
\label{sec:circuits_tp_light_cone}

\begin{definition}[Local properties and unbranched circuits]
    \label{def:circuits_local}
    \lean{QuantumCircuit.Circuit.Forall, QuantumCircuit.Circuit.IsUnbranched,
        QuantumCircuit.Circuit.lightCone}
    \leanok
    \uses{def:circuits_circuit}
    A property of Kraus families holds locally on a circuit when it holds locally on
    every layer of every branch. The backward light cone of a region $R$ through a
    circuit is obtained from the last layer to the first, each layer adding its
    cells that meet the current region, over all branches.
\end{definition}

\begin{lemma}[Local properties under composition]
    \label{lem:circuits_local_comp}
    \lean{QuantumCircuit.Circuit.Forall.append, QuantumCircuit.Circuit.Forall.comp,
        QuantumCircuit.Circuit.IsUnbranched.comp, QuantumCircuit.Layer.lightCone_mono,
        QuantumCircuit.Circuit.lightCone_mono, QuantumCircuit.Circuit.lightCone_comp}
    \leanok
    \uses{def:circuits_local, thm:circuits_composition}
    A property holding locally on a circuit and on every continuation attached to
    it holds locally on the composite. A composite of unbranched circuits is
    unbranched. The backward light cone is monotone in the region, and the backward
    light cone of $R$ through $c$ followed by $c'$ lies in the backward light cone
    through $c$ of the backward light cone of $R$ through $c'$.
\end{lemma}
\begin{proof}\leanok
    Induction on the first circuit; the layers of the composite are those of $c$
    followed by those of the continuations.
\end{proof}

\begin{theorem}[Trace preservation of a circuit]
    \label{thm:circuits_tp}
    \lean{QuantumCircuit.Circuit.dual_one_of_forall,
        QuantumCircuit.Circuit.isTP_krausFamily, QuantumCircuit.Circuit.isKrausCPTP_map,
        QuantumCircuit.Circuit.trace_map}
    \leanok
    \uses{def:circuits_local, def:kraus_family, def:is_kraus_cptp}
    A circuit whose gates and onsite families are trace-preserving is
    trace-preserving, hence a quantum channel, whether or not it branches.
\end{theorem}
\begin{proof}\leanok
    \uses{thm:circuits_layer_tp, lem:kraus_family_channel}
    By induction the Heisenberg picture fixes the identity. For a circuit that
    starts with a layer $L$ with Kraus operators $K_i$ and continues with $c_i$,
    \begin{align}
        \sum_iK_i^\dagger\Bigl(\sum_pK^{c_i\dagger}_p\,1\,K^{c_i}_p\Bigr)K_i
        &=\sum_iK_i^\dagger K_i=1,
        \notag
    \end{align}
    by the induction hypothesis in the middle and Theorem~\ref{thm:circuits_layer_tp}
    at the end. Then $\sum_pK_p^\dagger K_p=1$ for the branches $p$ of the
    circuit.
\end{proof}

\begin{theorem}[Backward light cone of a circuit]
    \label{thm:circuits_light_cone}
    \lean{QuantumCircuit.Circuit.dual_mem_supportedOperators_lightCone,
        QuantumCircuit.Layer.lightCone_subset_neighbourhood,
        QuantumCircuit.Circuit.lightCone_subset_neighbourhood}
    \leanok
    \uses{def:circuits_local, def:sites_supported}
    Let a circuit be unbranched with trace-preserving gates and onsite families, and
    let $A$ be supported on $R$. Then the Heisenberg picture of $A$ is supported on
    the backward light cone of $R$
    \cite{BravyiHastingsVerstraete2006LiebRobinson}. If the geometry has finitely
    many cells and every cell has positive cost, the backward light cone lies in the $T$-fold neighbourhood of $R$, where
    $T$ is the depth.
\end{theorem}
\begin{proof}\leanok
    \uses{thm:circuits_layer_light_cone, lem:circuits_neighbourhood_mono,
        lem:circuits_layer_factorization}
    Induction on the circuit. For an unbranched circuit the Heisenberg picture is
    that of the first layer applied to that of the common continuation, and
    Theorem~\ref{thm:circuits_layer_light_cone} applies. If every cell has positive
    cost, a layer of depth $0$ has no cells and adds nothing, and a layer of depth
    $1$ adds at most the neighbourhood. Both hypotheses are needed: with feedback
    from a measurement far from $R$ the Heisenberg picture depends on the measured
    site. Without trace preservation, in the factorization of
    Lemma~\ref{lem:circuits_layer_factorization} the far factor is the effect
    \begin{align}
        \sum_{i_F}(K^{\mathrm{far}}_{i_F})^\dagger K^{\mathrm{far}}_{i_F}&=E_F,
        \notag
    \end{align}
    so $\sum_iK_i^\dagger(A_1\otimes1)K_i
    =\bigl(\sum_{i_N}(K^{\mathrm{near}}_{i_N})^\dagger A_1K^{\mathrm{near}}_{i_N}\bigr)
    \otimes E_F$, which is not supported on $D$ when $E_F$ is not a multiple of the
    identity.
\end{proof}

\section{Sequences, layers of several gates and adjoints}
\label{sec:circuits_sequences}

\begin{lemma}[Sequences of layers of one operator]
    \label{lem:circuits_seq}
    \lean{QuantumCircuit.Circuit.kraus_seq, QuantumCircuit.Circuit.total_seq,
        QuantumCircuit.Circuit.forall_seq, QuantumCircuit.Circuit.isUnbranched_seq,
        QuantumCircuit.KrausFamily.isUnitary_single_one,
        QuantumCircuit.Layer.gateCount_single, QuantumCircuit.Layer.depth_single,
        QuantumCircuit.Layer.nonempty_index_single, QuantumCircuit.Layer.forall_single}
    \leanok
    \uses{def:circuits_circuit, def:circuits_costs, def:circuits_local,
        lem:circuits_layer_single}
    Let $L_1,\dots,L_m$ be layers on one system, and suppose every Kraus operator of
    $L_k$ equals $U_k$. Then the Kraus operator of the sequence along every branch is
    $U_m\cdots U_1$. The sequence is unbranched; a property holding locally on every
    layer holds locally on the sequence; and if every layer has at least one outcome,
    the cost of the sequence is the sum of the costs of its layers. The layer with a
    single gate on a cell $b$ has gate count $c_b$, and depth $1$ if $c_b>0$; the
    identity, as a family of one operator, is unitary.
\end{lemma}
\begin{proof}\leanok
    Induction on $m$. The Kraus operator of $L_1$ followed by the rest along a
    branch $(i,p)$ is the Kraus operator of the rest along $p$ times $U_1$. When
    every layer has an outcome, the largest cost of a continuation is the cost of
    the common continuation.
\end{proof}

\begin{definition}[A layer of gates on given cells]
    \label{def:circuits_of_cells}
    \lean{QuantumCircuit.Layer.ofCells, QuantumCircuit.Layer.cellOp,
        QuantumCircuit.Layer.placeFamily, QuantumCircuit.Layer.cellsUnion,
        QuantumCircuit.Layer.cellsProd}
    \leanok
    \uses{def:circuits_layer, def:sites_place}
    For a set of pairwise disjoint cells $b$ of a geometry and a gate $g_b$ on each,
    the layer of these gates carries the identity at every remaining site. For a
    plain gate $g$ on a region $D$ the family of placements of $g$ is the family of
    operators $K\otimes1$ on $\mathcal H_X$ for the operators $K$ of $g$.
\end{definition}

\begin{theorem}[A layer of disjoint gates is the product of their placements]
    \label{thm:circuits_of_cells}
    \lean{QuantumCircuit.Layer.kraus_ofCells, QuantumCircuit.Layer.kraus_ofCells_eq_prod,
        QuantumCircuit.Layer.map_ofCells, QuantumCircuit.Layer.cellsUnion_eq,
        QuantumCircuit.Layer.disjoint_cellsUnion,
        QuantumCircuit.Layer.prod_place_eq_place_cellsProd,
        QuantumCircuit.Layer.cellsProd_apply, QuantumCircuit.Layer.pairwiseDisjoint_of_cons,
        QuantumCircuit.Layer.consPi, QuantumCircuit.Layer.consIndexEquiv,
        QuantumCircuit.Layer.kraus_cons, QuantumCircuit.Layer.gateCount_ofCells,
        QuantumCircuit.Layer.depth_ofCells_le_one,
        QuantumCircuit.Layer.nonempty_index_ofCells}
    \leanok
    \uses{def:circuits_of_cells, def:circuits_costs, lem:sites_place_hom}
    Let $b_1,\dots,b_m$ be an enumeration without repetition of a set of pairwise
    disjoint cells carrying plain gates $g_{b_k}$. The Kraus operator of the layer at
    an outcome $i$ is
    \begin{align}
        K_i&=(K_{b_1,i_{b_1}}\otimes1)\cdots(K_{b_m,i_{b_m}}\otimes1),
        \notag
    \end{align}
    the product of the placements of the operators of the gates at $i$, and the map
    of the layer is the composite of the maps of the families of placements of the
    $g_{b_k}$, in this order. The gate count of the layer is the sum of the costs of
    its cells, and its depth is at most $1$.
\end{theorem}
\begin{proof}\leanok
    \uses{lem:circuits_layer_factorization}
    Placements on disjoint regions $D_1,D_2$ multiply to the placement on
    $D_1\cup D_2$ of the operator whose entry at $(y,x)$ is the product of the
    entries of the two factors at the restrictions of $y$ and $x$, and that entry is
    $0$ unless $y$ and $x$ agree outside $D_1\cup D_2$. By induction on $m$ the
    product of the placements is the placement on the union of the cells of the
    product of the local entries, which is the formula for the Kraus operator of a
    layer, the identity at the remaining sites contributing the condition that $y$
    and $x$ agree there. For the maps, split off the first cell: the outcomes of the
    layer are the pairs of an outcome of $g_{b_1}$ and an outcome of the layer of the
    other cells, and
    $\sum_{i,j}(K_iK'_j)\rho(K_iK'_j)^\dagger
    =\sum_iK_i\bigl(\sum_jK'_j\rho K_j'^\dagger\bigr)K_i^\dagger$.
\end{proof}

\begin{definition}[Adjoint of a layer and of a sequence]
    \label{def:circuits_adjoint}
    \lean{KrausFamily.adjoint, QuantumCircuit.Layer.adjoint,
        QuantumCircuit.Circuit.seqAdjoint}
    \leanok
    \uses{def:circuits_layer, def:kraus_family}
    The adjoint of a finite Kraus family $(K_i)_i$ is $(K_i^\dagger)_i$. The adjoint
    of a layer of plain gates from $X$ to $Y$ is the layer from $Y$ to $X$ on the
    same cells, with the adjoint gates and onsite families. The adjoint of a sequence
    of layers of plain gates on one system is the sequence of the adjoint layers in
    reverse order.
\end{definition}

\begin{lemma}[Kraus operators and costs of adjoints]
    \label{lem:circuits_adjoint}
    \lean{QuantumCircuit.Layer.kraus_adjoint, QuantumCircuit.Layer.adjoint_adjoint,
        QuantumCircuit.Circuit.kraus_seqAdjoint, QuantumCircuit.Circuit.total_seqAdjoint}
    \leanok
    \uses{def:circuits_adjoint, lem:circuits_seq}
    The Kraus operators of the adjoint layer are the adjoints of those of the layer,
    and the adjoint of the adjoint layer is the layer. If every Kraus operator of the
    $k$-th layer of a sequence equals $U_k$, the Kraus operator of the adjoint
    sequence is $(U_m\cdots U_1)^\dagger$. A cost unchanged by taking the adjoint
    layer, such as depth or gate count, is the same for a sequence of layers with
    outcomes and for its adjoint.
\end{lemma}
\begin{proof}\leanok
    The entry of a Kraus operator of a layer is a product of entries of the gates and
    onsite families; conjugating it and exchanging the configurations gives the
    entry of the adjoint. The rest follows from Lemma~\ref{lem:circuits_seq}, with
    $U_1^\dagger\cdots U_m^\dagger=(U_m\cdots U_1)^\dagger$ and the sum of costs
    unchanged by reversing the order.
\end{proof}

An adjoint is defined only for a sequence on one system: when the systems change
from layer to layer the Kraus operators need not be square, and when later layers
depend on outcomes the reversal of a branch is not a circuit of the same kind.

\section{Gates on sites and pairs}
\label{sec:circuits_local_gates}

\begin{definition}[Gates from matrices on a local basis]
    \label{def:circuits_local_gate}
    \lean{QuantumCircuit.siteEquiv, QuantumCircuit.pairEquiv,
        QuantumCircuit.pairEquiv_symm_apply_left, QuantumCircuit.pairEquiv_symm_apply_right,
        QuantumCircuit.KrausGate.ofMatrix, QuantumCircuit.Geometry.sitesAndPairs}
    \leanok
    \uses{def:circuits_gate_family, def:circuits_geometry, def:sites_config}
    The configurations of a single site $a$ are the basis $X_a$, and those of
    $\{a,b\}$ for $a\neq b$ are the pairs in $X_a\times X_b$, the first factor at
    $a$. A matrix on a set $B$ in bijection with the configurations of a region $D$
    gives the plain gate on $D$ with this one Kraus operator. The geometry of sites
    and pairs on $V$ has a cell $\{a\}$ of cost $c_1$ for every site and a cell
    $\{a,b\}$ of cost $c_2$ for every ordered pair of distinct sites.
\end{definition}

\begin{lemma}[A unitary matrix gives a unitary gate]
    \label{lem:circuits_local_gate_unitary}
    \lean{QuantumCircuit.KrausGate.isUnitary_ofMatrix}
    \leanok
    \uses{def:circuits_local_gate}
    The gate of a unitary matrix on a local basis is unitary.
\end{lemma}
\begin{proof}\leanok
    Relabelling rows and columns along one bijection preserves unitarity.
\end{proof}

\section{Moving circuits}
\label{sec:circuits_maps}

A circuit moves to a coarser geometry or a larger gate set on the same sites, and to
a larger set of sites in which its sites are embedded. In both cases its Kraus
operators, local properties, light cones and costs move with it.

\begin{definition}[Morphisms of geometries and of gate families]
    \label{def:circuits_geometry_hom}
    \lean{QuantumCircuit.Geometry.Hom, QuantumCircuit.Geometry.Hom.cell_eq,
        QuantumCircuit.Geometry.Hom.cellsEquiv, QuantumCircuit.GateMorphism}
    \leanok
    \uses{def:circuits_geometry, def:circuits_gate_family}
    Let $\iota:V\to V'$ be injective. A morphism of geometries over $\iota$ is an
    injective map of cells $b\mapsto b'$ with $C_{b'}=\iota(C_b)$. A morphism of gate
    families on the same sites, compatible with interpretations, maps every gate on
    $D$ from $X$ to $Y$ to a gate on $D$ from $X$ to $Y$, with a bijection of
    outcomes under which the interpreted matrices agree.
\end{definition}

\begin{theorem}[Moving a circuit on the same sites]
    \label{thm:circuits_map_gate}
    \lean{QuantumCircuit.Layer.mapGate, QuantumCircuit.Layer.mapGateCellsEquiv,
        QuantumCircuit.Layer.mapGateFreeEquiv, QuantumCircuit.Layer.mapGateIndexEquiv,
        QuantumCircuit.Layer.kraus_mapGate, QuantumCircuit.Layer.lightCone_mapGate,
        QuantumCircuit.Layer.gateCount_mapGate, QuantumCircuit.Layer.depth_mapGate_le,
        QuantumCircuit.Circuit.mapGate, QuantumCircuit.Circuit.mapGateIndexEquiv,
        QuantumCircuit.Circuit.kraus_mapGate, QuantumCircuit.Circuit.isUnbranched_mapGate,
        QuantumCircuit.Circuit.forall_mapGate, QuantumCircuit.Circuit.total_mapGate_le,
        QuantumCircuit.Circuit.gateCount_mapGate_le, QuantumCircuit.Circuit.depth_mapGate_le,
        QuantumCircuit.Circuit.lightCone_mapGate}
    \leanok
    \uses{def:circuits_geometry_hom, def:circuits_circuit, def:circuits_costs,
        def:circuits_local}
    Along a morphism of geometries over the identity of $V$ and a compatible morphism
    of gate families, every circuit has an image on the same systems, whose branches
    are in bijection with those of the circuit and have the same Kraus operators. The
    image is unbranched if the circuit is, a property of Kraus families invariant
    under relabelling of outcomes holds locally on the image if it holds locally on
    the circuit, and the light cones agree. If no cell costs more after the move, the
    gate count does not grow; if no cell of cost $0$ gets positive cost, the depth
    does not grow.
\end{theorem}
\begin{proof}\leanok
    The cells of the image layer are the images of the cells of the layer, with the
    same regions, so the remaining sites and the onsite families do not change. The
    outcomes correspond through the bijections of outcomes of the gates, and the
    product of local entries is the same product. The statements about circuits
    follow by induction.
\end{proof}

\begin{definition}[Systems on a larger set of sites]
    \label{def:circuits_lift}
    \lean{QuantumCircuit.siteEmb, QuantumCircuit.siteEmb_apply, QuantumCircuit.lift,
        QuantumCircuit.Idle, QuantumCircuit.siteSplit, QuantumCircuit.siteEmb_siteSplit,
        QuantumCircuit.liftSiteEquiv, QuantumCircuit.liftConfigEquiv,
        QuantumCircuit.liftConfigEquiv_fst, QuantumCircuit.liftConfigEquiv_snd,
        QuantumCircuit.regionSiteEquiv, QuantumCircuit.liftRegionEquiv}
    \leanok
    \uses{def:sites_config}
    Let $V'$ be in bijection with the disjoint union of $V$ and a finite set $R$ of
    idle sites, each $r\in R$ with a basis $W_r$, and let $\iota:V\to V'$ be the
    inclusion of $V$. A system $X$ on $V$ becomes the system on $V'$ with basis
    $X_s$ at $\iota(s)$ and $W_r$ at $r$; its configurations are the pairs of a
    configuration of $X$ and a configuration of the idle sites, and the
    configurations of $\iota(D)$ are those of $D$. Every injective map $\iota$ arises
    this way, with $R$ the complement of its image.
\end{definition}

\begin{definition}[Transport of gates]
    \label{def:circuits_transport}
    \lean{QuantumCircuit.GateTransport, QuantumCircuit.GateTransport.IsCompatible,
        QuantumCircuit.KrausGate.transport, QuantumCircuit.KrausGate.isCompatible_transport}
    \leanok
    \uses{def:circuits_lift, def:circuits_gate_family}
    A transport of gate families along $V'\cong V\sqcup R$ maps every gate on $D$
    from $X$ to $Y$ to a gate on $\iota(D)$ between the systems on $V'$, with a
    bijection of outcomes. It is compatible with interpretations when the
    interpreted matrices agree under the identification of configurations of $D$
    and of $\iota(D)$. Plain gates have a compatible transport, the same operators
    read on $\iota(D)$.
\end{definition}

\begin{theorem}[Moving a circuit to a larger set of sites]
    \label{thm:circuits_map_sites}
    \lean{QuantumCircuit.Layer.cellsEquiv, QuantumCircuit.Layer.cell_cellsEquiv_symm,
        QuantumCircuit.Layer.liftOnsite, QuantumCircuit.Layer.not_mem_of_mapSites,
        QuantumCircuit.Layer.mapSites, QuantumCircuit.Layer.mapSitesCellsEquiv,
        QuantumCircuit.Layer.freeEquiv, QuantumCircuit.Layer.mapSitesIndexEquiv,
        QuantumCircuit.Layer.kraus_mapSites, QuantumCircuit.Layer.lightCone_mapSites,
        QuantumCircuit.Layer.gateCount_mapSites, QuantumCircuit.Layer.depth_mapSites_le,
        QuantumCircuit.Circuit.mapSites, QuantumCircuit.Circuit.mapSitesIndexEquiv,
        QuantumCircuit.Circuit.kraus_mapSites, QuantumCircuit.Circuit.isUnbranched_mapSites,
        QuantumCircuit.Circuit.total_mapSites_le, QuantumCircuit.Circuit.forall_mapSites,
        QuantumCircuit.Circuit.lightCone_mapSites,
        QuantumCircuit.Circuit.gateCount_mapSites_le,
        QuantumCircuit.Circuit.depth_mapSites_le}
    \leanok
    \uses{def:circuits_lift, def:circuits_transport, def:circuits_geometry_hom,
        def:circuits_circuit, def:circuits_costs, def:circuits_local}
    Along a morphism of geometries over $\iota$ and a compatible transport of gates,
    every circuit from $X$ to $Z$ on $V$ has an image from the lifted $X$ to the
    lifted $Z$ on $V'$, carrying the identity at the idle sites. Its branches are in
    bijection with those of the circuit, and along corresponding branches
    \begin{align}
        K'_p&=K_p\otimes1_{W},
        \notag
    \end{align}
    in the identification of configurations of $V'$ with pairs. The image is
    unbranched if the circuit is; a property of Kraus families invariant under
    relabelling of bases and outcomes, and holding for the identity, holds locally
    on the image if it holds locally on the circuit; the light cone of $\iota(T)$ is
    the image of the light cone of $T$; and the gate count and the depth do not grow
    under the conditions of Theorem~\ref{thm:circuits_map_gate}.
\end{theorem}
\begin{proof}\leanok
    The cells of the image layer are the images of the cells, and its remaining sites
    are the images of the remaining sites of the layer together with the idle sites.
    The outcomes are those of the gates, of the onsite families at the images, and
    the single outcome of the identity at each idle site. In the product of local
    entries the factors at the cells and at the images of the remaining sites are
    those of the layer, and the factors at the idle sites form the identity on
    $\C^{W}$. The statements about circuits follow by induction, a branch of the
    image being a branch of the circuit read through these bijections.
\end{proof}
